% Part: model-theory
% Chapter: interpolation
% Section: introduction

\documentclass[../../../include/open-logic-section]{subfiles}

\begin{document}

\olfileid[pt]{mod}{int}{int}

\olsection{Introdução}

O teorema da interpolação é o seguinte resultado: suponha que
$\Entails !A \lif !B$. Então existe !!a{sentence} $!C$ tal que
$\Entails !A \lif !C$ e $\Entails !C \lif !B$. Além disso, cada
!!{constant}, !!{function} e !!{predicate} (diferente de $\eq$) que ocorre em
$!C$ ocorre tanto em $!A$ quanto em~$!B$. A !!{sentence} $!C$ é chamada de
\emph{interpolante} de $!A$ e~$!B$.

O teorema da interpolação é interessante por si só, mas sua principal
importância está no fato de poder ser usado para demonstrar resultados sobre
definibilidade em uma teoria e sobre as condições nas quais a combinação de
duas teorias consistentes resulta em uma teoria consistente. O primeiro
resultado é conhecido como teorema da definibilidade de Beth; o segundo, como
teorema da consistência conjunta de Robinson.

\end{document}
